\documentclass[11pt,a4paper]{article}
\usepackage{fontspec,amsmath,amssymb,amsthm,microtype}
\setmainfont{lmroman10-regular.otf}[BoldFont=lmroman10-bold.otf,ItalicFont=lmroman10-italic.otf,BoldItalicFont=lmroman10-bolditalic.otf]
\usepackage[margin=25mm]{geometry}
\usepackage[unicode,hidelinks]{hyperref}
\hypersetup{pdfauthor={LangevinLean Project}}
\setlength{\emergencystretch}{3em}
\setlength{\parskip}{0.4em}
\setlength{\parindent}{0pt}
\allowdisplaybreaks[2]
\raggedbottom
\title{\texorpdfstring{Interacting Particle Langevin under Strong Convexity}{Interacting Particle Langevin under Strong Convexity}}
\author{LangevinLean Project}
\date{}
\begin{document}\maketitle
Draft mathematical exposition; these notes may contain mistakes.
The bound follows from strong convexity and gradient Lipschitzness of the joint potential. Parameter and latent dimensions need not agree. Normalization, moments and concentration of the Gibbs measure are proved under these assumptions. The separate entropy result assumes LSI for the joint particle target and finite initial joint relative entropy.
\section{\texorpdfstring{Theorem}{Theorem}}


Fix an observation \(y\) and a differentiable joint potential \(U_y:\mathbb R^{d_\theta}\times\mathbb R^{d_x}\to\mathbb R\). Assume \(U_y\) is \(m\)-strongly convex with \(m>0\) and has \(L\)-Lipschitz gradient. Let \(N\ge1\); the independent dimensions \(d_\theta,d_x\ge0\) may be zero. Define


\[\ell_y(\theta)=\int e^{-U_y(\theta,x)}dx,\quad
\theta_\star\in\operatorname*{argmax}_\theta\ell_y(\theta).\]


In scaled coordinates \(z=(\theta,x^1/\sqrt N,\ldots,x^N/\sqrt N)\), set


\[V_N(z)=\sum_{i=1}^NU_y(\theta,\sqrt N z_i),\qquad
\Pi_N(dz)\propto e^{-V_N(z)}dz.\]


Assume the initial state has a finite second moment. Let all Gaussian vectors in the following recursion be standard and independent of one another and of the initial state. For \(0<\gamma\le2/(m+L)\), define the IPLA iterates by


\[\begin{aligned}
\theta_{k+1}&=\theta_k-\frac\gamma N\sum_i\nabla_\theta U_y(\theta_k,x_k^i)+\sqrt{2\gamma/N}\,\xi_k^0,\\
x_{k+1}^i&=x_k^i-\gamma\nabla_xU_y(\theta_k,x_k^i)+\sqrt{2\gamma}\,\xi_k^i.
\end{aligned}\]


Writing \(\mu_0\) for the initial scaled-state law, for every \(k\in\mathbb N\),


\[\begin{aligned}
\bigl(\mathbb E\|\theta_k-\theta_\star\|^2\bigr)^{1/2}
\le{}&\sqrt{\frac{2d_\theta}{Nm}}+(1-m\gamma)^kW_2(\mu_0,\Pi_N)\\
&+\frac{33}{20}\frac Lm\sqrt{d_x+\frac{d_\theta}{N}}\sqrt\gamma.
\end{aligned}\]


The assumptions imply that the Gibbs measures are normalized by finite positive constants and have finite second moments. The proof also establishes strong convexity of the negative log marginal. Existence and uniqueness of its minimizer are proved separately from the same assumptions, so \(\theta_\star\) exists and is the unique maximizer of the integrated joint density \(\ell_y\). If the potential includes the negative log of a prior density, the same definition maximizes that prior-weighted marginal. The proof compares the particle distribution with its Gibbs measure and then bounds the parameter's second moment about \(\theta_\star\).


\section{\texorpdfstring{The scaled particle geometry}{The scaled particle geometry}}



For \(N\ge1\) use coordinates \(z=(\theta,x^1/\sqrt N,\ldots,x^N/\sqrt N)\) with squared norm \(\|\theta\|^2+N^{-1}\sum_i\|x^i\|^2\). The pair map \(J_i z=(\theta,x^i)\) satisfies

\[\sum_{i=1}^N\|J_i(z-z')\|^2=N\|z-z'\|^2.\]

Summing the pair strong-convexity inequalities gives quadratic correction \(Nm\|z-z'\|^2/2\). Thus \(V_N=\sum_iU_y\circ J_i\) has curvature \(Nm\). For the gradient constant \(NL\), apply the chain rule through \(J_i\), test a gradient difference in an arbitrary direction, and apply Cauchy–Schwarz over particles. The same squared-norm identity bounds both factors, yielding the stated constant.



The dimension is \(D=d_\theta+Nd_x\). The argument permits unequal parameter and latent dimensions, and zero-dimensional blocks contribute zero to the sums.



\section{\texorpdfstring{The particle recursion as a Gaussian Euler update}{The particle recursion as a Gaussian Euler update}}



The chain rule in the parameter and latent blocks gives

\[\nabla_\theta V_N(z)=\sum_i\nabla_\theta U_y(\theta,x^i),\qquad
\nabla_{z_i}V_N(z)=\sqrt N\nabla_xU_y(\theta,x^i).\]

Use Gaussian Euler for \(V_N\) at step \(s=\gamma/N\). The parameter drift is the average pair gradient and its Gaussian coefficient is \(\sqrt{2\gamma/N}\). Multiplying a scaled latent update by \(\sqrt N\) gives drift coefficient \(\gamma\) and noise coefficient \(\sqrt{2\gamma}\). These are the physical IPLA updates.



The independent next innovation is independent of the current state, so its conditional update is the Gaussian Euler transition. Iterating this transition gives the particle laws. The restriction \(s\le2/(Nm+NL)\) is equivalent to \(\gamma\le2/(m+L)\) under positive particle count.



\section{\texorpdfstring{Latent integration and marginal strong convexity}{Latent integration and marginal strong convexity}}



The normalization and second-moment theorems \cite[Section~3.26.13]{LangevinLean2026ULASC} and \cite[Section~3.26.14]{LangevinLean2026ULASC} apply to the strongly convex summed potential. Change the scaled latent blocks back to physical vectors. Its volume factor is independent of \(\theta\) and cancels on normalization. Tonelli gives

\[\int\exp\!\left(-\sum_iU_y(\theta,x^i)\right)\prod_i dx^i
=\prod_i\int e^{-U_y(\theta,x)}dx=\ell_y(\theta)^N.\]

The nonnegative product integral factors before finiteness is used. Gaussian domination and positivity then make \(F(\theta)=-\log\ell_y(\theta)\) an ordinary real marginal potential.



To prove curvature of \(F\), subtract the parameter quadratic corresponding to \(m\) from \(U_y\). The remaining exponential is jointly log-concave. The finite-dimensional Prékopa–Leindler theorem preserves log-concavity under latent integration, and restoration of the quadratic gives the claimed strong convexity. Continuity and coercivity give a minimizer, whose uniqueness follows from strict convexity. It is the marginal-likelihood maximizer \(\theta_\star\).



\section{\texorpdfstring{Equilibrium fluctuations about the marginal maximizer}{Equilibrium fluctuations about the marginal maximizer}}



The parameter marginal potential \(NF\) has curvature \(Nm\) and minimizer \(\theta_\star\). The continuous strongly convex concentration theorem \cite[Section~3.49.7]{LangevinLean2026ULASC} gives

\[\int\|\theta-\theta_\star\|^2\,\Pi_N(dz)\le\frac{2d_\theta}{Nm}.\]

Its nonnegative square root is \(r_N=\sqrt{2d_\theta/(Nm)}\). This is the radius about the marginal mode; minimizing over latent coordinates instead of integrating them would generally give a different point. Although the dynamics have dimension \(D\), only parameter dimension appears in this concentration estimate. Under \(d_\theta=0\) the projection is constant and the radius vanishes. A separate dilation argument establishes the improved factor used in later proximal and kinetic estimates.



\section{\texorpdfstring{The parameter projection of a transport coupling}{The parameter projection of a transport coupling}}



Define \(Pz=\theta\). Orthogonality of the state blocks gives \(\|P(z-z')\|\le\|z-z'\|\). For a coupling \((Z,Z_*)\) of \(\mu_k\) and \(\Pi_N\), Minkowski yields

\[\begin{aligned}
\|PZ-\theta_\star\|_{L^2}
&\le\|P(Z-Z_*)\|_{L^2}+\|PZ_*-\theta_\star\|_{L^2}\\
&\le\|Z-Z_*\|_{L^2}+r_N.
\end{aligned}\]

Taking the finite-cost coupling infimum, as in \cite[Section~3.48.9]{LangevinLean2026ULASC}, gives \(\|\theta_k-\theta_\star\|_{L^2}\le W_2(\mu_k,\Pi_N)+r_N\). Finite second moments of both laws justify these real quantities. No particle independence at equilibrium or finite time is needed for this projection argument.



\section{\texorpdfstring{The particle-scaled transport and parameter bound}{The particle-scaled transport and parameter bound}}



Use curvature \(Nm\), gradient constant \(NL\), dimension \(D=d_\theta+Nd_x\), and step \(s=\gamma/N\) in the Gaussian Euler estimate \cite[Section~3.45.7]{LangevinLean2026ULASC}. The contraction becomes \(1-(Nm)s=1-m\gamma\) and the dimension term satisfies

\[\frac{NL}{Nm}\sqrt{sD}
=\frac Lm\sqrt{\gamma\left(d_x+\frac{d_\theta}{N}\right)}.\]

Adding the equilibrium radius by the preceding projection inequality proves the parameter estimate. The powers of \(N\) are determined by the coordinate and time scalings. Parameter noise and equilibrium spread decrease with \(N\), while the latent discretization contribution has order \(\sqrt{\gamma d_x}\).



Thus \(N\) controls equilibrium concentration, \(\gamma\) controls the time-discretization error, and \(k\) controls the remaining initial error. At a fixed physical step particle count alone does not remove the latent numerical error. Unequal dimensions are included by the same identities.



\begin{thebibliography}{ULA-LSI}
\bibitem[ULA-SC]{LangevinLean2026ULASC}
LangevinLean Project. \emph{Unadjusted Langevin under Strong Convexity}. LL-01, 2026. Draft report, content revision c58ab5c70c3a. \url{https://odakyildiz.github.io/LangevinLean/docs/pdf/UnadjustedLangevin/FullProof.pdf}.
\end{thebibliography}

\end{document}
